\documentclass[10pt]{article}

\usepackage[utf8]{inputenc}

\usepackage[T1]{fontenc}

\usepackage{amsmath}

\usepackage{amsfonts}

\usepackage{amssymb}

\usepackage[version=4]{mhchem}

\usepackage{stmaryrd}

\usepackage{fixltx2e}



\title{62 БРАГИНСКОГО }



\author{}

\date{}





\begin{document}

\maketitle

Сходимость Б. р. определяется свойствами потенциала $v(x)$. Для любых короткодействующих потенциалов Б. р. сходится при достаточно больших $|k|$. Если же потенциал $v(x)$ удовлетворяет условию



\[

\int|v(x)||v(y)||x-y|^{-2} d x d y<4 \pi

\]



то Б. р. сходится при всех значениях параметра $k$. При $|x| \rightarrow \infty$ формула (1) дает Б. р. для амплитуды рассеяния $f^{( \pm)}(\hat{x}, k)$ :



\[

\begin{gathered}

\Psi^{( \pm)}(x, k)=\Psi_{0}(x, k)+f^{( \pm)}(\hat{x}, k)\left[\frac{\exp \{i|k||x|\}}{|x|}+o(1)\right], \\

f^{( \pm)}(\hat{x}, k)=\sum_{n=0}^{\infty} f_{n}^{( \pm)}(\hat{x}, k), \\

f_{n}^{( \pm)}(\hat{x}, k)=-(4 \pi)^{-1} \int \exp \{\mp i|k|(\hat{x}, k)\} v(y) \Psi_{n}^{( \pm)}(y, k) d y .

\end{gathered}

\]



Первое слагаемое в (2)



\[

f_{0}^{( \pm)}(\hat{x}, k)=-(4 \pi)^{-1} \int \exp \{\mp i|k|(\hat{x}-k, y)\} v(y) d y

\]



называется борновским приближением дляамплитуды рассеяния.



Jum.: [1] B orn M., «Z. Phys.», 1926, Bd 38, S. 803; [2] M отт H., Мес с и Г., Теория атомных столкновений, пер. с англ., 3 изд., М., 1969; [3] Рид М., С ай мон Б., Методы современной математической физики, пер. с англ., т. 3. Теория рассеяния, М., 1982.



С. П. Меркурьев, С.Л. Яковлев.



БОРНОВСКОЕ ПРИБЛИЖЕНИЕ в к в а то во Й М е ханике и квантовой те ори и поля-приближенный метод вычисления амплитуд упругого рассеяния и неупругого взаимодействия микрочастиц в рамках возмущений теории в первом приближении по потенциалу взаимодействия. Метод сформулирован М. Борном (M. Born, 1926). Применимость Б. п. для короткодействующих потенциалов определяется условием $U R / \hbar v \ll 1$, где $R-$ размер области действия потенциала, $v$ - относительная скорость рассеиваемых частиц, $\bar{U}$ - среднее значение потенциала (в случае квантовой теории поля - энергии взаимодействия) в области с размером $\sim R$. Это условие означает, что время $\sim R / v$, к-рое частицы проводят в области взаимодействия, мало по сравнению со временем $\sim \hbar / U$, за к-рое взаимодействие успевает сильно изменить состояние частиц. Для кулоновского поля Б. П. справедливо при условии $Z \alpha / \hbar v \ll 1$, где $Z$ - атомный номер, $\alpha \approx 1 / 137$ - постоянная тонкой структуры. Это означает, что скорость $v$ частиц должна превышать скорость $v_{\mathrm{K}}=Z \alpha / \hbar$ движения электрона на первой боровской орбите. Б. П. лучше выполняется при больших скоростях частиц. При произвольных $v$ оно справедливо, если $|U(R)| \ll \hbar^{2} / m R^{2}$.



В нерелятивистской квантовой механике при справедливости Б. п. амплитуда упругого рассеяния действительна и равна



\[

f=-\frac{m}{2 \pi \hbar^{2}} \int U(\boldsymbol{r}) e^{-i k \boldsymbol{r}} d V

\]



где $\hbar \boldsymbol{k}=\boldsymbol{p}_{\boldsymbol{i}}-\boldsymbol{p}_{f}-$ изменение импульса в процессе рассеяния, $\boldsymbol{p}_{i}$ и $\boldsymbol{p}_{f}-$ импульсы рассеиваемых частиц до и после рассеяния, $m$ - масса рассеиваемой частицы, $U(\boldsymbol{r})$ - потенциал взаимодействия ( $d V$ - элемент объема).



Поскольку в общем случае амплитуда рассеяния является комплексной величиной, ее действительность в Б. п. означает, что фазы рассеяния $\delta_{l}$ в состоянии с орбитальным квантовым числом $l$ должны быть малы. Для них в Б. П. справедливо выражение:



\[

\delta_{l}=-\frac{\pi m}{2 \hbar} \int_{0}^{\infty} U(\boldsymbol{r})\left[J_{l+1 / 2}(k r)\right]^{2} \boldsymbol{r} d \boldsymbol{r}

\]



где $J_{l+1 / 2}$ - функция Бесселя.



Б. п. широко используется при анализе упругого и неупругого рассеяния и служит основным методом извлече- ния информации о формфакторах элементарных частиц, атомов и атомных ядер.



См. также Борновский ряд. Л. И. Лапидус, М. В. Терентьев. БОУЭНА-МАРКУСА-МАРГУЛИСА МЕРА - см. Орициклический поток.



БOХHEPA ИНTETPAЛ - интеграл от функции со значениями в банаховом пространстве по скалярной мере. Пусть $\left\{x_{n}(s)\right\}$ - последовательность функций на измеримом пространстве $(S, \mathfrak{B}, m)$, принимающих конечное число значений, со значением в банаховом пространстве, сильно сходящаяся к функции $x(s)$ почти всюду. Й т т грал Бохн е ра функции $x(s)$ по мере $m(d s)$ определен как сильный предел последовательности интегралов $x_{n}(s)$ по мере $m(d s)$, то есть



\[

\int_{S} x(s) m(d s)=s t \lim _{n \rightarrow \infty} \int_{S} x_{n}(s) m(d s) .

\]



В. И. Скрипник.



БРАГИНСКОГО УРАВНЕНИЯ - система Гидродинамических уравнений для полностью ионизованной плазмы, состоящей из электронов и одного сорта ионов с зарядом $e$. Они состоят из уравнен и й не преры вности:



\[

\frac{\partial n^{(\alpha)}}{\partial t}+\operatorname{div}\left(n^{(\alpha)} \boldsymbol{V}^{(\alpha)}\right)=0,

\]



уравнений движения:



\[

\begin{aligned}

& m^{(\alpha)} n^{(\alpha)} \frac{\partial^{(\alpha)} V_{p}^{(\alpha)}}{\partial t}=-\frac{\partial p^{(\alpha)}}{\partial x_{p}}-\frac{\partial \pi_{p q}^{(\alpha)}}{\partial x_{q}}+ \\

& +e^{(\alpha)} n^{(\alpha)}\left\{E_{p}+\frac{1}{c}\left(V^{(\alpha)} \times B\right)_{p}\right\}+R_{p}^{(\alpha)}

\end{aligned}

\]



и уравнений теплового баланса:



\[

\frac{3}{2} n^{(\alpha)} \frac{\partial^{(\alpha)} T^{(\alpha)}}{\partial t}+p^{(\alpha)} \operatorname{div} V^{(\alpha)}=-\operatorname{div} \boldsymbol{q}^{(\alpha)}-\pi_{p q}^{(\alpha)} \frac{\partial V_{p}^{(\alpha)}}{\partial x_{q}}+Q^{(\alpha)} .

\]



Здесь индекс $\alpha$ означает сорт частиц $(\alpha=e$ для электронов и $\alpha=i$ для ионов), $n^{(\alpha)}-$ плотность, $V^{(\alpha)}-$ скорость, $m^{(\alpha)}-$ масса, $\frac{d^{(\alpha)}}{d t}=\frac{\partial}{\partial t}+\left(V^{(\alpha)} \nabla\right), p^{(\alpha)}=n^{(\alpha)} T^{(\alpha)}-$ давление, $T^{(\alpha)}-$ температура в энергетич. единицах, $e^{(e)}=-e$, $e^{(i)}=e, \boldsymbol{R}^{(\alpha)}-$ передача импульса из-за столкновений с частицами другого сорта,



\[

\boldsymbol{R}^{(e)}=-\boldsymbol{R}^{(i)}=\boldsymbol{R}^{j}+\boldsymbol{R}^{T} \equiv \boldsymbol{R},

\]



$\boldsymbol{R}^{j}$ - сила трения, $\boldsymbol{R}^{T}-$ термосила,



\[

\begin{gathered}

\boldsymbol{R}^{j}=e n\left(\frac{j_{\|}}{\sigma_{\|}}+\frac{j_{\perp}}{\sigma_{\perp}}\right), \\

\boldsymbol{R}^{T}=-0,71 n \nabla_{\|} T^{(e)}-\frac{3}{2} \frac{n}{\omega^{(e)_{\tau}(e)}} \frac{\boldsymbol{B}}{B} \times \nabla T^{(e)} .

\end{gathered}

\]



Значки | и $\perp$ означают параллельное и перпендикулярное вектору направления, $\omega^{(\alpha)}=\frac{e^{(\alpha)} B}{m^{(\alpha)} c}, \frac{1}{\tau^{(\alpha)}}-$ частота столкновений частиц сорта $\alpha$ друг с другом, $\left|\omega^{(\alpha)}\right| \tau^{(\alpha)} \gg 1$. Во всех соотношениях, кроме уравнения Максвелла



\[

\operatorname{div} \boldsymbol{E}=4 \pi e\left(n^{(i)}-n^{(e)}\right),

\]



предполагается выполненным условие квазинейтральности



\[

n^{(e)}=n^{(i)} \equiv n

\]



$\boldsymbol{q}^{(\alpha)}-$ плотность потока тепла:



\[

\boldsymbol{q}^{(\alpha)}=\boldsymbol{q}^{(\alpha) j}+\boldsymbol{q}^{(\alpha) T}

\][



\textbackslash boldsymbol\{q\}\^{}\{(e) j\}=-0,71 \textbackslash frac\{T\^{}\{(e)\}\}\{e\} \textbackslash boldsymbol\{j\}\_\{|\}-\textbackslash frac\{3\}\{2\} \textbackslash frac\{T\^{}\{(e)\}\}\{e \textbackslash omega\^{}\{(e)\} \textbackslash tau\^{}\{(e)\}\} \textbackslash frac\{B\}\{B\} \textbackslash times \textbackslash boldsymbol\{j\}, \textbackslash quad \textbackslash boldsymbol\{q\}\^{}\{(i) j\} \textbackslash approx 0,



][



\textbackslash boldsymbol\{q\}\^{}\{(\textbackslash alpha) T\}=-x\_\{|\}\^{}\{(\textbackslash alpha)\} \textbackslash nabla\_\{|\} T\textsuperscript{\{(\textbackslash alpha)\}-x\_\{\textbackslash perp\}}\{(\textbackslash alpha)\} \textbackslash nabla\_\{\textbackslash perp\} T\^{}\{(\textbackslash alpha)\}+\textbackslash frac\{5\}\{2\} \textbackslash frac\{n T\textsuperscript{\{(\textbackslash alpha)\}\}\{\textbackslash omega}\{(\textbackslash alpha)\} m\^{}\{(\textbackslash alpha)\}\} \textbackslash frac\{B\}\{B\} \textbackslash times \textbackslash nabla T\^{}\{(\textbackslash alpha)\},

]



$Q^{(\alpha)}-$ тепло, получаемое частицами сорта $\alpha$ вследствие столкновений с частицами другого сорта,



\[

\begin{gathered}

Q^{(i)}=\frac{3 m^{(e)}}{m^{(i)}} \frac{n}{\tau^{(e)}}\left(T^{(e)}-T^{(i)}\right), \\

Q^{(e)}=\frac{1}{e n} \boldsymbol{R j}-Q^{(i)} .

\end{gathered}

\]





\end{document}